% Generated draft: requires mathematical and source review.
\documentclass[11pt]{article}
\usepackage[margin=25mm]{geometry}
\usepackage{fontspec}
\setmainfont{Latin Modern Roman}
\usepackage{amsmath,amssymb,mathtools,unicode-math}
\setmathfont{Latin Modern Math}
\usepackage{enumitem,needspace,xurl,hyperref}
\setlength{\parindent}{0pt}
\setlength{\parskip}{5pt}
\setlength{\emergencystretch}{3em}
\newcommand{\N}{\mathbb{N}}
\newcommand{\Z}{\mathbb{Z}}
\newcommand{\Q}{\mathbb{Q}}
\newcommand{\R}{\mathbb{R}}
\newcommand{\C}{\mathbb{C}}
\newcommand{\F}{\mathbb{F}}
\newcommand{\A}{\mathbb{A}}
\newcommand{\PP}{\mathbb{P}}
\newcommand{\E}{\mathbb{E}}
\newcommand{\eps}{\varepsilon}
\newcommand{\dd}{\,\mathrm{d}}
\begin{document}
\section*{UnsolvedMath ID: 30002397}
\textbf{Source problem code:} catalogue entry
\subsection{Definitions and mathematical statement}
The supplied record asks a question about a method of proof, but does not specify enough hypotheses or an analytic target to determine a unique mathematical conjecture.

One conventional framework in which to interpret its terminology is the following; these conventions are explanatory and are not additional hypotheses established by the source record. Let \((M,\tau)\) be a tracial von Neumann probability space, and let \(X_1,\ldots,X_n\in M\) be selfadjoint elements. Their generated unital algebras are freely independent if
\[
\tau(a_1\cdots a_k)=0
\]
whenever each \(a_j\) belongs to one of these algebras, \(\tau(a_j)=0\), and consecutive factors belong to different algebras. For a selfadjoint \(Y\in M\), its scalar distribution is the probability measure
\[
\mu_Y(B)=\tau\bigl(\mathbf{1}_B(Y)\bigr),\qquad B\subseteq\R\text{ Borel}.
\]
It is atomless if \(\mu_Y(\{t\})=0\) for every \(t\in\R\). For a noncommutative polynomial \(p=p^*\in\C\langle x_1,\ldots,x_n\rangle\), where the formal variables are selfadjoint, write
\[
Y=p(X_1,\ldots,X_n).
\]
Under the interpretation that ``free variables without atoms'' refers to the individual input distributions, the input condition would be
\[
\mu_{X_i}(\{t\})=0\qquad(1\leq i\leq n,\ t\in\R).
\]
The source does not explicitly specify this framework, whether the variables are bounded or selfadjoint, or whether this is the intended meaning of ``without atoms.''

For terminology, given a conditional expectation \(E_{\mathcal B}:\mathcal A\to\mathcal B\) between unital operator algebras and a selfadjoint \(L\in\mathcal A\), an operator-valued Cauchy transform has the form
\[
G_L(b)=E_{\mathcal B}\bigl[(b-L)^{-1}\bigr]
\]
for admissible \(b\in\mathcal B\) for which the inverse exists. A linearization description of a polynomial uses an auxiliary matrix pencil, affine in the input variables, whose resolvent encodes the polynomial resolvent. The supplied record does not identify the algebras, conditional expectation, pencil, resolvent relation, or analytic domain intended here.

\textbf{Source question.} For a selfadjoint polynomial in free variables without atoms, can the absence of atoms in its distribution, and similar qualitative results, be derived from the operator-valued Cauchy-transform and linearization description?

With the conventional interpretation above, the named distributional target would be
\[
\mu_{p(X_1,\ldots,X_n)}(\{t\})=0\qquad\text{for every }t\in\R.
\]
This display identifies the requested conclusion, not a universally quantified assertion justified by the supplied formulation. In particular, the source gives no restriction excluding constant polynomials, and ``similar qualitative results'' specifies no further properties. The permitted polynomial class, precise input hypotheses, and meaning of deriving the result from the proposed description therefore remain unresolved.
\subsection{Short English statement}
The question asks whether operator-valued Cauchy transforms and polynomial linearization can yield atomlessness of the distribution of a selfadjoint polynomial in freely independent atomless inputs, together with other unspecified qualitative properties. The input setting, polynomial restrictions, analytic description, and additional properties need clarification before this becomes a precise conjecture. The supplied record labels it open; this draft does not independently verify that status.
\subsection{Sources}
\begin{itemize}
\item Local UnsolvedMath problem registry, permanent ID \(30002397\), title ``Atomlessness from Operator-Valued Cauchy Transforms,'' source problem code \texttt{catalogue entry}, in the Oberwolfach Reports Open Problems collection. No individual proposer is supplied; the proposed year is recorded as 2013.
\item Supplied report citation: \emph{C*-Algebren (2014)}, Oberwolfach Reports, pp.\ 2423--2500, DOI \texttt{10.4171/owr/2013/43}. Supplied source URL: \url{https://doi.org/10.4171/owr/2013/43}. The background describes Oberwolfach Reports as records of mathematical workshops and their open-problem discussions. No problem-specific page or named contributor is supplied.
\item Supplied literature source: ``Oberwolfach Report 43,'' \url{https://doi.org/10.4171/owr/2013/43}. Literature checked at: \texttt{2026-08-21}. This is supplied dated curation, not independently verified. Its status label is ``open,'' and its assessment says that the extracted statement is too vague to identify a precise theorem or conjecture and that the cited report does not establish a resolution. The background requests independent verification of time-sensitive status and cited claims.
\item The registry marks the statement as ``exact'' and supplies a verification note claiming that its mathematical content was checked against the official report, with notation, context, and question form normalized without changing the source-level claim. That verification is part of the supplied provenance and has not been independently repeated here. The citation's 2014 label and the recorded 2013 proposed year and DOI year are retained as supplied.
\end{itemize}
\end{document}
